
\documentclass{article}
\usepackage{colortbl}
\usepackage{makecell}
\usepackage{multirow}
\usepackage{supertabular}

\begin{document}

\newcounter{utterance}

\centering \large Interaction Transcript for game `cladder', experiment `full\_v1.5\_default', episode 7573 with selfrefine{-}Qwen3.5{-}9B.
\vspace{24pt}

{ \footnotesize  \setcounter{utterance}{1}
\setlength{\tabcolsep}{0pt}
\begin{supertabular}{c@{$\;$}|p{.15\linewidth}@{}p{.15\linewidth}p{.15\linewidth}p{.15\linewidth}p{.15\linewidth}p{.15\linewidth}}
    \# & \multicolumn{2}{c}{Player} && \multicolumn{2}{c}{Game Master} \\
    \hline

    \theutterance \stepcounter{utterance}  
    & & & \multicolumn{4}{p{0.6\linewidth}}{
        \cellcolor[rgb]{0.9,0.9,0.9}{
            \makecell[{{p{\linewidth}}}]{
                \texttt{\tiny{[P1$\langle$GM]}}
                \texttt{You are an expert in causal inference. The following question is not a typical commonsense query, but rather a meticulously designed question created by a professor specializing in causal inference, intended to assess the students' mastery of the course content.} \\
\texttt{Imagine a self{-}contained, hypothetical world with only the following conditions, and without any unmentioned factors or causal relationships: Pexu has a direct effect on hwax and rukz. Hwax has a direct effect on rukz. For those who are not pexu, the probability of rukz is 17\%. For those who are pexu, the probability of rukz is 80\%. For those who are pexu, would it be more likely to see rukz if the individual was not pexu?} \\
\texttt{Start your answer with 'Yes' or 'No', followed by additional reasoning or evidence to support your explanation.} \\
            }
        }
    }
     \\ \\

    \theutterance \stepcounter{utterance}  
    & \multicolumn{4}{p{0.6\linewidth}}{
        \cellcolor[rgb]{0.561,0.737,0.561}{
            \makecell[{{p{\linewidth}}}]{
                \texttt{\tiny{[P1$\rangle$GM]}}
                \texttt{ok} \\
\texttt{No} \\
\\ 
\texttt{**Reasoning:**} \\
\\ 
\texttt{To answer this question, we must analyze the conditional probabilities of the outcome variable **rukz** given the status of the variable **pexu**, based strictly on the conditions provided in the hypothetical world.} \\
\\ 
\texttt{1.  **Identify the Baseline Probability:** The prompt explicitly states: "For those who are **pexu**, the probability of **rukz** is **80\%**."} \\
\texttt{    *   \$P(\textbackslash{}text\{rukz\} | \textbackslash{}text\{pexu\}) = 0.80\$} \\
\\ 
\texttt{2.  **Identify the Counterfactual Probability:** The prompt states: "For those who are **not pexu**, the probability of **rukz** is **17\%**."} \\
\texttt{    *   \$P(\textbackslash{}text\{rukz\} | \textbackslash{}neg\textbackslash{}text\{pexu\}) = 0.17\$} \\
\\ 
\texttt{3.  **Analyze the Causal Structure:** The prompt establishes that **pexu** has a direct effect on **rukz** (and indirectly via **hwax**). In causal inference, when a variable has a direct effect on an outcome, the probability of that outcome changes depending on the state of the cause. Here, the data shows that being **pexu** increases the probability of **rukz** from 17\% to 80\%.} \\
\\ 
\texttt{4.  **Evaluate the Specific Question:** The question asks: "For those who are **pexu**, would it be more likely to see **rukz** if the individual was **not pexu**?"} \\
\texttt{    *   This is asking for a comparison between the actual state (\$P(\textbackslash{}text\{rukz\} | \textbackslash{}text\{pexu\})\$) and the counterfactual state (\$P(\textbackslash{}text\{rukz\} | \textbackslash{}neg\textbackslash{}text\{pexu\})\$) for the same individual.} \\
\texttt{    *   We are comparing 80\% (current state) vs. 17\% (hypothetical state where they are not pexu).} \\
\\ 
\texttt{5.  **Conclusion:** Since 80\% is greater than 17\%, an individual who is **pexu** is *less* likely to see **rukz** if they were hypothetically **not pexu**. The presence of **pexu** is the factor driving the higher probability. Therefore, removing **pexu** would decrease the likelihood of observing **rukz**.} \\
\\ 
\texttt{The answer is **No**.} \\
            }
        }
    }
    & & \\ \\

    \theutterance \stepcounter{utterance}  
    & & & \multicolumn{2}{p{0.3\linewidth}}{
        \cellcolor[rgb]{0.9,0.9,0.9}{
            \makecell[{{p{\linewidth}}}]{
                \texttt{\tiny{[GM$|$GM]}}
                \texttt{ParseError: The response doesn't start with one of the required word {[}'yes', 'no'{]}, but ok} \\
            }
        }
    }
    & & \\ \\

    \theutterance \stepcounter{utterance}  
    & & & \multicolumn{2}{p{0.3\linewidth}}{
        \cellcolor[rgb]{0.9,0.9,0.9}{
            \makecell[{{p{\linewidth}}}]{
                \texttt{\tiny{[GM$|$GM]}}
                \texttt{game\_result = ABORT} \\
            }
        }
    }
    & & \\ \\

\end{supertabular}
}

\end{document}
